% Chapitre 1 : une histoire de commits, avec une tag, une branche et sa fusion.
\begin{figure}[tbp]
\centering
\begin{tikzpicture}[x=1cm,y=1cm,
  commit/.style={circle, draw=black!60, fill=white, line width=0.8pt, minimum size=2.6mm, inner sep=0pt},
  commitb/.style={commit, draw=tonregle},
  hash/.style={font=\ttfamily\scriptsize, text=black!55, anchor=west},
  msg/.style={font=\sffamily\scriptsize, anchor=west}]
  % les deux lignes
  \draw[black!60, line width=1.2pt] (0,0) -- (0,4.5);
  \draw[tonregle, line width=1.2pt] (0,1.5) to[out=90,in=-90] (0.9,2.25) -- (0.9,3.75)
    to[out=90,in=-90] (0,4.5);
  % les commits, du plus ancien (en bas) au plus récent
  \foreach \y/\x/\s/\h/\m in {
      0/0/commit/3f9a1c2/Premier jet de l'article,
      0.75/0/commit/b81e07d/Introduction : le plan et les notations,
      1.5/0/commit/4233568/Lemme 3 : la finitude est nécessaire (cf. 3.4),
      2.25/0.9/commitb/c0d5e11/Lemme 5 : l'énoncé{,} preuve à faire,
      3/0/commit/7a4b9f0/Section 2 : une coquille dans la définition 2.1,
      3.75/0.9/commitb/e92f6a8/Lemme 5 : preuve complète{,} axiomes vérifiés,
      4.5/0/commit/1d7c3b5/Fusion de la pull request \#12 (lemme 5)} {
    \node[\s] at (\x,\y) {};
    \node[hash] at (2.5,\y) {\h};
    \node[msg] at (3.85,\y) {\m};
  }
  % la tag
  \node[etiquette, draw=tonregle, fill=tonregle!8, text=tonregle, anchor=east]
    at (-0.3,1.5) {\ttfamily article-soumis-2026-10};
  \draw[tonregle, densely dotted] (-0.3,1.5) -- (-0.13,1.5);
  % les noms de branches
  \node[etiq, anchor=south] at (0,4.65) {\ttfamily main};
  \node[etiq, text=tonregle, anchor=west] at (1.0,2.62) {\ttfamily lemme-5};
  \node[etiq, text=tonregle, anchor=west] at (0.75,4.12) {pull request};
\end{tikzpicture}
\caption{Une histoire de commits, du plus ancien en bas au plus récent en haut. Chaque
commit porte son identifiant abrégé et son message. Le tag nomme l'état envoyé à
l'éditeur ; le lemme~5 a été écrit sur une branche, puis fusionné dans \texttt{main} par
une pull request relue.}
\label{fig:historique}
\end{figure}
